\section{Moving the analytic class to the two ends}\label{sec:sliding}

The two endpoint maps $i_-f,i_+f:Y\to X$ are connected componentwise
by the height slices of $X_j$.  The distance between the endpoints grows
with $j$, so a single bounded-distance homotopy does not compare their
Roe-algebra maps.  Instead, we place a copy of an operator at every
height and compare this stack with its shift by one height.  The common
interior copies will cancel on $K$-theory.

\begin{proposition}\label{ana:endpoints}
The endpoint maps induce the same homomorphism
\[
 (i_{+,A}f_A)_*=(i_{-,A}f_A)_*:
 K_1(A_Y)\longrightarrow K_1(A_X).
\]
\end{proposition}

For $-3j\le t\le3j$, let
\[
 s_{j,t}:Y_j\longrightarrow X_j,
 \qquad (u,q)\longmapsto (u,\overline q,t),
\]
where $\overline q$ is reduction from $Q_{2j}$ to $Q_{k_j(t)}$.
The map is $1$-Lipschitz: it carries both Cayley edges and horizontal
matchings to edges or to vertices.  Consecutive slices satisfy
\begin{equation}\label{eq:successive-slices}
 d_X(s_{j,t}(y),s_{j,t+1}(y))=1.
\end{equation}
Choose exact covering isometries
\[
 D_{j,t}:H_{Y_j}\longrightarrow H_{X_j}
\]
over these vertex maps.  They exist because every vertex summand is a
copy of the infinite-dimensional Hilbert space $H_0$.  At the endpoints
choose them to be the restrictions of the covering isometries defining
$i_{-,A}f_A$ and $i_{+,A}f_A$, respectively.  For fixed $j$, the ranges
of the $D_{j,t}$ for distinct $t$ are orthogonal, since they lie at
distinct heights.

\subsection{Stacks modulo the compact operators}

Let
\[
 q:A_X\longrightarrow A_X/\K_X
\]
be the quotient map.  If $T\in\C[Y]$, write
$T_j=\chi_{Y_j}T\chi_{Y_j}$.  The difference
$T-\bigoplus_jT_j$ has finite vertex support: at a fixed propagation
bound there are only finitely many cross-component entries.  In
particular it is compact.  Define
\begin{align}
 S_0(T)&=\bigoplus_j\sum_{t=-3j}^{3j-1}
                D_{j,t}T_jD_{j,t}^*,\label{eq:stack-zero}\\
 S_1(T)&=\bigoplus_j\sum_{t=-3j}^{3j-1}
                D_{j,t+1}T_jD_{j,t+1}^*.
                \label{eq:stack-one}
\end{align}
The number of copies is unbounded, but their ranges are orthogonal.
Thus $\|S_a(T)\|_{\mathrm{op}}\le\|T\|_{\mathrm{op}}$, for $a=0,1$.
If $\prop(T)\le R$, the $1$-Lipschitz property gives
$\prop(S_a(T))\le R$.
Every vertex entry is compact: it is a finite sum of isometric images
of compact entries of $T_j$, since each $Y_j$ is finite.  Hence
$S_a(T)\in\C[X]$.  No bound on the ranks of these compact entries is
required in the algebra used here.

The maps $S_a$ need not be multiplicative before taking the quotient.
For $T,T'\in\C[Y]$, however,
\begin{equation}\label{eq:compression-defect}
 (TT')_j-T_jT'_j
   =\chi_{Y_j}T(1-\chi_{Y_j})T'\chi_{Y_j}
\end{equation}
is zero except for finitely many $j$.  Each nonzero term is compact on
$H_{Y_j}$, since it is a matrix of compact operators over a finite
vertex set.  Copying it into the finitely many heights of the
corresponding $X_j$ remains compact.  Orthogonality of the height
ranges and \eqref{eq:compression-defect} therefore show that
\[
 \overline S_a=qS_a:\C[Y]\longrightarrow A_X/\K_X
\]
is a $*$-homomorphism.  By maximal universality, it extends to a
$*$-homomorphism with domain $A_Y$.

\subsection{A continuous rotation between the stacks}

We next prove that $\overline S_0$ and $\overline S_1$ induce the same
map on $K$-theory.  The rotation takes place in a doubled target module;
this ensures that the ranges stay orthogonal at every intermediate
angle.

For $0\le\lambda\le\pi/2$, define
\[
 W_{\lambda,j,t}=
 \begin{pmatrix}
  \cos\lambda\,D_{j,t}\\
  \sin\lambda\,D_{j,t+1}
 \end{pmatrix}
 :H_{Y_j}\longrightarrow H_{X_j}\oplus H_{X_j},
 \qquad -3j\le t<3j.
\]
Each $W_{\lambda,j,t}$ is an isometry.  For distinct pairs $(j,t)$,
their ranges are orthogonal: they lie in different components or, within
one component, at different heights in each target copy.  In the first
copy the heights are $t$, and in the second they are $t+1$.
Thus they assemble to an
isometry
\[
 W_\lambda:\bigoplus_{j}\bigoplus_{t=-3j}^{3j-1}H_{Y_j}
          \longrightarrow H_X\oplus H_X.
\]
On $\C[Y]$ define
\begin{equation}\label{eq:stack-rotation}
 F_\lambda(T)=W_\lambda
       \left(\bigoplus_j\bigoplus_{t=-3j}^{3j-1}T_j\right)
       W_\lambda^*.
\end{equation}
Its diagonal matrix blocks are scalar multiples of the stacks in
\eqref{eq:stack-zero}--\eqref{eq:stack-one}.  The off-diagonal blocks
are sums of $D_{j,t}T_jD_{j,t+1}^*$ and their counterparts.  If
$\prop(T)\le R$, then \eqref{eq:successive-slices} gives propagation
at most $R+1$ for these blocks.  Their vertex entries are again finite
sums of compact operators.  We conclude that
\[
 F_\lambda(T)\in M_2(\C[X]),\qquad
 \prop(F_\lambda(T))\le R+1.
\]
The same calculation as \eqref{eq:compression-defect}, using
$W_\lambda^*W_\lambda=1$, shows that the images of $F_\lambda$ in
$M_2(A_X/\K_X)$ are $*$-homomorphisms on $\C[Y]$.  Each extends to
$A_Y$ by maximal universality.

For completeness, continuity is uniform in the number of height
slices.  The orthogonality just established gives
\[
 \|W_\lambda-W_\mu\|
  =\bigl((\cos\lambda-\cos\mu)^2
            +(\sin\lambda-\sin\mu)^2\bigr)^{1/2}
  \le |\lambda-\mu|.
\]
Consequently, in concrete operator norm,
\begin{equation}\label{eq:rotation-continuity}
 \|F_\lambda(T)-F_\mu(T)\|_{\mathrm{op}}
       \le2\|T\|_{\mathrm{op}}\,|\lambda-\mu|.
\end{equation}
The difference has propagation at most $R+1$, independently of the
angles.  Applying Lemma~\ref{ana:fixed-propagation} to its four matrix
blocks promotes \eqref{eq:rotation-continuity} to an estimate
\[
 \|F_\lambda(T)-F_\mu(T)\|_{M_2(A_X)}
       \le C_R\|T\|_{\mathrm{op}}\,|\lambda-\mu|
\]
for a constant independent of $j,\lambda,\mu$.
This proves point-norm continuity on the dense algebra $\C[Y]$.
The extended homomorphisms are contractive for the maximal norm, so
density proves point-norm continuity on all of $A_Y$.

At $\lambda=0$ the homomorphism is $\overline S_0$ in the first
matrix corner, and at $\lambda=\pi/2$ it is $\overline S_1$ in the
second.  The corner identifications on $K$-theory agree, so
\begin{equation}\label{eq:equal-stacks}
 (\overline S_0)_*=(\overline S_1)_*.
\end{equation}

\subsection{Canceling the interior copies}

\begin{proof}[Proof of Proposition~\ref{ana:endpoints}]
Define the interior stack by
\[
 I(T)=\bigoplus_j\sum_{t=-3j+1}^{3j-1}D_{j,t}T_jD_{j,t}^*.
\]
Its image $\overline I=qI$ is a $*$-homomorphism on $A_Y$ by the
same compact-defect argument.  Let
\[
 E_-=i_{-,A}f_A,\qquad E_+=i_{+,A}f_A.
\]
The endpoints in the two stacks agree modulo $\K_X$ with $E_-$ and
$E_+$, respectively.  Indeed, our endpoint covering isometries agree
with those defining $E_\pm$, and discarding the cross-component part
of $T$ changes its image by a compact operator.  Hence
\[
 \overline S_0=\overline I+qE_-,\qquad
 \overline S_1=\overline I+qE_+.
\]
In each sum, the two homomorphisms have orthogonal ranges: one is
supported on the strictly interior heights and the other on an
endpoint.

Orthogonal sums induce sums on $K_1$.  Explicitly, if $\varphi$ and
$\psi$ are orthogonal homomorphisms and $u$ is a matrix unitary with
$u-1$ in the domain algebra, then
\[
 1+(\varphi+\psi)(u-1)
   =\bigl(1+\varphi(u-1)\bigr)
      \bigl(1+\psi(u-1)\bigr),
\]
and multiplication of unitaries gives addition of their $K_1$-classes.
Using this additivity in \eqref{eq:equal-stacks} and canceling
$\overline I_*$ gives
\[
 q_*(E_-)_*=q_*(E_+)_*.
\]
By Lemma~\ref{ana:compacts}, $\K_X$ is the algebra of compact
operators on $H_X$, so $K_1(\K_X)=0$.  The six-term exact sequence
for
\[
 0\longrightarrow\K_X\longrightarrow A_X
       \xrightarrow{q}A_X/\K_X\longrightarrow0
\]
therefore makes $q_*:K_1(A_X)\to K_1(A_X/\K_X)$ injective.
It follows that $(E_-)_*=(E_+)_*$, as asserted.
\end{proof}

\begin{proof}[Proof of Theorem~\ref{thm:main}]
The graphs $X_j$ constructed above are finite and connected with a
uniform degree bound, and their coarse disjoint union $X$ is uniformly
discrete and has bounded geometry.  Proposition~\ref{top:infinite-order}
gives an infinite-order class
\[
 \alpha=(i_+)_*\beta-(i_-)_*\beta\in\KX_1(X).
\]
By Proposition~\ref{ana:transfer}, choose $b\in K_1(A_Y)$ satisfying
\[
 (f_A)_*(b)=\mu_B^{\max}(\beta).
\]
Naturality of the maximal coarse index
(Proposition~\ref{ana:naturality}) and
Proposition~\ref{ana:endpoints} now give
\begin{align*}
 \mu_X^{\max}(\alpha)
  &=(i_{+,A})_*\mu_B^{\max}(\beta)
       -(i_{-,A})_*\mu_B^{\max}(\beta)\\
  &=(i_{+,A}f_A)_*(b)-(i_{-,A}f_A)_*(b)=0.
\end{align*}
An infinite-order element remains nonzero after tensoring with
$\mathbb Q$, so $\alpha\otimes1$ is a nonzero element of the kernel
of the rationalized maximal coarse assembly map.
\end{proof}
